% Generated mechanically from the frozen stacks-perfect.tex target.
% Do not edit this generated file; edit the bound locale source instead.

% OK, start here.
%

\setcounter{chapter}{103}
\chapter{叠的导出范畴}


\phantomsection
\label{stacks-perfect-section-phantom}





\section{引言}
\label{stacks-perfect-section-introduction}

\noindent
本章讨论与代数叠相伴的导出范畴，尤其包括拟凝聚层的导出范畴；
也就是说，我们将证明概形情形（参见《概形的导出范畴》第
\ref{perfect-section-introduction} 节）和代数空间情形
（参见《空间的导出范畴》第
\ref{spaces-perfect-section-introduction} 节）相应结果的类似结论。
本章结果与 \cite{LM-B} 中结果的主要区别是，我们始终使用“大 site”。
阅读本章之前，请先快速浏览《代数叠上的层》和《代数叠的上同调》两章，
其中引入了本章所用的术语。



\section{约定、记号与语言滥用}
\label{stacks-perfect-section-conventions}

\noindent
我们继续采用《叠的性质》第
\ref{stacks-properties-section-conventions} 节中引入的约定和语言滥用。
记号依照《叠的上同调》第
\ref{stacks-cohomology-section-notation} 节。












\section{光滑-\'etale site 与平坦-fppf site}
\label{stacks-perfect-section-lisse-etale}

\noindent
本节是《叠的上同调》第
\ref{stacks-cohomology-section-lisse-etale} 节在导出范畴情形下的类似版本。

\begin{lemma}
\label{stacks-perfect-lemma-shriek-derived}
设 $\mathcal{X}$ 为代数叠。采用《叠的上同调》引理
\ref{stacks-cohomology-lemma-lisse-etale} 和
\ref{stacks-cohomology-lemma-lisse-etale-modules} 中的记号。
\begin{enumerate}
\item 函子
$g_! : \textit{Ab}(\mathcal{X}_{lisse,\etale}) \to
\textit{Ab}(\mathcal{X}_\etale)$
有左导出函子
$$
Lg_! :
D(\mathcal{X}_{lisse,\etale})
\longrightarrow
D(\mathcal{X}_\etale)
$$
它是 $g^{-1}$ 的左伴随，并且 $g^{-1}Lg_! = \text{id}$。
\item 函子 $g_! : 
\textit{Mod}(\mathcal{X}_{lisse,\etale},
\mathcal{O}_{\mathcal{X}_{lisse,\etale}}) \to
\textit{Mod}(\mathcal{X}_\etale, \mathcal{O}_{\mathcal{X}})$
有左导出函子
$$
Lg_! :
D(\mathcal{O}_{\mathcal{X}_{lisse,\etale}})
\longrightarrow
D(\mathcal{X}_\etale, \mathcal{O}_\mathcal{X})
$$
它是 $g^*$ 的左伴随，并且 $g^*Lg_! = \text{id}$。
\item 函子 $g_! : \textit{Ab}(\mathcal{X}_{flat,fppf}) \to
\textit{Ab}(\mathcal{X}_{fppf})$
有左导出函子
$$
Lg_! :
D(\mathcal{X}_{flat, fppf})
\longrightarrow
D(\mathcal{X}_{fppf})
$$
它是 $g^{-1}$ 的左伴随，并且 $g^{-1}Lg_! = \text{id}$。
\item 函子 $g_! :
\textit{Mod}(\mathcal{X}_{flat,fppf},
\mathcal{O}_{\mathcal{X}_{flat,fppf}}) \to
\textit{Mod}(\mathcal{X}_{fppf}, \mathcal{O}_{\mathcal{X}})$
有左导出函子
$$
Lg_! :
D(\mathcal{O}_{\mathcal{X}_{flat, fppf}})
\longrightarrow
D(\mathcal{O}_\mathcal{X})
$$
它是 $g^*$ 的左伴随，并且 $g^*Lg_! = \text{id}$。
\end{enumerate}
警告：先验地并不清楚模上的 $Lg_!$ 是否与阿贝尔层上的 $Lg_!$ 一致；
参见《site 上的上同调》评注
\ref{sites-cohomology-remark-when-derived-shriek-equal}。
\end{lemma}

\begin{proof}
函子 $Lg_!$ 的存在性及其与 $g^*$ 的伴随性由《site 上的上同调》引理
\ref{sites-cohomology-lemma-existence-derived-lower-shriek} 给出。
（对阿贝尔层的情形，取常值层 $\mathbf{Z}$ 为结构层。）
此外，它在复形 $\mathcal{H}^\bullet$ 上的计算方式是：
取一个适当的左分解
$\mathcal{K}^\bullet \to \mathcal{H}^\bullet$，
再将函子 $g_!$ 作用于 $\mathcal{K}^\bullet$。
由《叠的上同调》引理
\ref{stacks-cohomology-lemma-lisse-etale-modules} 和
\ref{stacks-cohomology-lemma-lisse-etale}，
$g^{-1}g_!\mathcal{K}^\bullet = \mathcal{K}^\bullet$，
故每种情形下的最后一个断言都成立。
\end{proof}










\begin{lemma}
\label{stacks-perfect-lemma-lisse-etale-functorial-derived}
采用《叠的上同调》引理
\ref{stacks-cohomology-lemma-lisse-etale-functorial} 中的假设和记号。
在无界导出范畴上有
$$
g^{-1} \circ Rf_* = Rf'_* \circ (g')^{-1}
\quad\text{且}\quad
L(g')_! \circ (f')^{-1} = f^{-1} \circ Lg_!
$$
（模的情形和阿贝尔层的情形均成立）。
\end{lemma}

\begin{proof}
令 $\tau = \etale$（相应地，令 $\tau = fppf$）。
设 $\mathcal{F}$ 为 $\mathcal{X}_\tau$ 上的阿贝尔层。
由《叠的上同调》引理
\ref{stacks-cohomology-lemma-lisse-etale-functorial-cohomology}，
典范的（基变换）映射
$$
g^{-1}Rf_*\mathcal{F} \longrightarrow Rf'_*(g')^{-1}\mathcal{F}
$$
是同构。证明的其余部分是形式的。由于阿贝尔群层与模层的上同调一致，
当 $\mathcal{F}$ 是 $\mathcal{X}_\tau$ 上的模层时，我们也得到
$g^{-1} Rf_*\mathcal{F} = Rf'_* (g')^{-1}\mathcal{F}$。

\medskip\noindent
接下来证明：对 $\mathcal{Y}_{lisse,\etale}$（相应地，
$\mathcal{Y}_{flat,fppf}$）上的 $\mathcal{G}$
（模层或阿贝尔群层），典范映射
$$
L(g')_!(f')^{-1}\mathcal{G} \to f^{-1}Lg_!\mathcal{G}
$$
是同构。为此，只需证明：对 $\mathcal{X}_\tau$ 上任意内射层
$\mathcal{I}$，诱导映射
$$
\Hom(L(g')_!(f')^{-1}\mathcal{G}, \mathcal{I}[n])
\leftarrow
\Hom(f^{-1}Lg_!\mathcal{G}, \mathcal{I}[n])
$$
对所有 $n \in \mathbf{Z}$ 都是同构。（这里的 Hom 取自适当的导出范畴。）
由 $f^{-1}$ 与 $Rf_*$ 的伴随性、$Lg_!$ 与 $g^{-1}$ 的伴随性，
以及相应的“加撇”版本，这归结为上面已经证明的同构
$g^{-1} Rf_*\mathcal{I} \to Rf'_* (g')^{-1}\mathcal{I}$。

\medskip\noindent
% P12-ERR-0042: see ../control/ERRATA_CANDIDATES.jsonl
若 $\mathcal{G}^\bullet$ 是
$\mathcal{Y}_{lisse,\etale}$（相应地，$\mathcal{Y}_{fppf}$）
上的有界复形（模复形或阿贝尔群复形），则典范映射
\begin{equation}
\label{stacks-perfect-equation-to-show}
L(g')_!(f')^{-1}\mathcal{G}^\bullet \to f^{-1}Lg_!\mathcal{G}^\bullet
\end{equation}
是同构。这由层的情形、涉及截断的通常论证，以及
$L(g')_!(f')^{-1}$ 和 $f^{-1}Lg_!$ 都是三角范畴的正合函子这一事实得到。

\medskip\noindent
设 $\mathcal{G}^\bullet$ 是
$\mathcal{Y}_{lisse,\etale}$（相应地，$\mathcal{Y}_{fppf}$）
上的上有界复形（模复形或阿贝尔群复形）。
典范映射（\ref{stacks-perfect-equation-to-show}）是同构，因为可用笨截断
$\sigma_{\geq -n}$（参见《同调》第
\ref{homology-section-truncations} 节）将
$\mathcal{G}^\bullet$ 写成有界复形的余极限
$\mathcal{G}^\bullet = \colim \mathcal{G}_n^\bullet$。
由此得到优三角
$$
\bigoplus\nolimits_{n \geq 1} \mathcal{G}_n^\bullet \to
\bigoplus\nolimits_{n \geq 1} \mathcal{G}_n^\bullet \to
\mathcal{G}^\bullet \to \ldots
$$
而函子 $L(g')_!$、$(f')^{-1}$、$f^{-1}$、$Lg_!$
中的每一个都与（复形的）直和交换。

\medskip\noindent
若 $\mathcal{G}^\bullet$ 是
$\mathcal{Y}_{lisse,\etale}$（相应地，$\mathcal{Y}_{fppf}$）
上的任意复形（模复形或阿贝尔群复形），则用典范截断
$\tau_{\leq n}$（参见《同调》第
\ref{homology-section-truncations} 节）将
$\mathcal{G}^\bullet$ 写成上有界复形的余极限，并重复上一段的论证。

\medskip\noindent
最后，由 $f^{-1}$ 与 $Rf_*$ 的伴随性、$Lg_!$ 与 $g^{-1}$ 的伴随性，
以及相应的“加撇”版本，我们得到：在完全一般的情形下，
引理的第一个恒等式由第二个恒等式推出。
\end{proof}
\begin{lemma}
\label{stacks-perfect-lemma-higher-shriek-quasi-coherent}
设 $\mathcal{X}$ 为代数叠。采用《叠的上同调》引理
\ref{stacks-cohomology-lemma-lisse-etale} 中的记号。
\begin{enumerate}
\item 设 $\mathcal{H}$ 为 $\mathcal{X}$ 的光滑-\'etale site 上的拟凝聚
$\mathcal{O}_{\mathcal{X}_{lisse,\etale}}$-模。
对所有 $p \in \mathbf{Z}$，层 $H^p(Lg_!\mathcal{H})$
都是 $\mathcal{X}$ 上具有平坦基变换性质的局部拟凝聚模。
\item 设 $\mathcal{H}$ 为 $\mathcal{X}$ 的平坦-fppf site 上的拟凝聚
$\mathcal{O}_{\mathcal{X}_{flat,fppf}}$-模。
对所有 $p \in \mathbf{Z}$，层 $H^p(Lg_!\mathcal{H})$
都是 $\mathcal{X}$ 上具有平坦基变换性质的局部拟凝聚模。
\end{enumerate}
\end{lemma}

\begin{proof}
选择概形 $U$ 以及满的光滑态射 $x : U \to \mathcal{X}$。
由《site 上的模》定义
\ref{sites-modules-definition-site-local}，存在一个 \'etale
（相应地，fppf）覆盖 $\{U_i \to U\}_{i \in I}$，
使得每个拉回 $f_i^{-1}\mathcal{H}$ 都有全局表示
（参见《site 上的模》定义
\ref{sites-modules-definition-global}）。
这里 $f_i : U_i \to \mathcal{X}$ 是复合
$U_i \to U \to \mathcal{X}$，因而是代数叠之间的态射。
（回忆拉回“就是”在 $\mathcal{X}/f_i$ 上的限制；参见《叠上的层》定义
\ref{stacks-sheaves-definition-pullback} 及其后的讨论。）
加细覆盖后，可设每个 $U_i$ 都是仿射概形。
由于每个 $f_i$ 都是光滑的（相应地，平坦的），由引理
\ref{stacks-perfect-lemma-lisse-etale-functorial-derived} 可知
$f_i^{-1}Lg_!\mathcal{H} = Lg_{i, !}(f'_i)^{-1}\mathcal{H}$。
利用《叠的上同调》引理
\ref{stacks-cohomology-lemma-check-lqc-fbc-on-covering}，
我们将引理的断言约化到如下情形：
$\mathcal{H}$ 有全局表示，并且对某个仿射概形
$X = \Spec(A)$ 有 $\mathcal{X} = (\Sch/X)_{fppf}$。

\medskip\noindent
设我们的表示形如
$$
\bigoplus\nolimits_{j \in J} \mathcal{O} \longrightarrow
\bigoplus\nolimits_{i \in I} \mathcal{O} \longrightarrow
\mathcal{H} \longrightarrow 0
$$
其中 $\mathcal{O} = \mathcal{O}_{\mathcal{X}_{lisse,\etale}}$
（相应地，$\mathcal{O} = \mathcal{O}_{\mathcal{X}_{flat,fppf}}$）。
注意 site $\mathcal{X}_{lisse,\etale}$
（相应地，$\mathcal{X}_{flat,fppf}$）有终对象，即
$X/X$，而且它是拟紧的（参见《site 上的上同调》第
\ref{sites-cohomology-section-limits} 节）。
因此由《Site》引理
\ref{sites-lemma-directed-colimits-sections} 有
$$
\Gamma(\bigoplus\nolimits_{i \in I} \mathcal{O}) =
\bigoplus\nolimits_{i \in I} A
$$
所以表示中的映射对应于 $A$-模 $M$ 的类似表示
$$
\bigoplus\nolimits_{j \in J} A \longrightarrow
\bigoplus\nolimits_{i \in I} A \longrightarrow
M \longrightarrow 0
$$
此外，$\mathcal{H}$ 等于与 $M$ 相伴的拟凝聚层 $M^a$
在光滑-\'etale（相应地，平坦-fppf）site 上的限制。
选择一个由自由 $A$-模组成的分解
$$
\ldots \to F_2 \to F_1 \to F_0 \to M \to 0
$$
复形
$$
\ldots \mathcal{O} \otimes_A F_2 \to \mathcal{O} \otimes_A F_1 \to
\mathcal{O} \otimes_A F_0 \to \mathcal{H} \to 0
$$
是 $\mathcal{H}$ 的自由 $\mathcal{O}$-模分解，因为对
$\mathcal{X}_{lisse,\etale}$（相应地，
$\mathcal{X}_{flat,fppf}$）的每个对象 $U/X$，
结构态射 $U \to X$ 都是平坦的。
因此按构造，$Lg_!\mathcal{H}$ 的值为
$$
\ldots \to
\mathcal{O}_\mathcal{X} \otimes_A F_2 \to
\mathcal{O}_\mathcal{X} \otimes_A F_1 \to
\mathcal{O}_\mathcal{X} \otimes_A F_0 \to 0 \to \ldots
$$
由于这是 $\mathcal{X}_\etale$（相应地，
$\mathcal{X}_{fppf}$）上的拟凝聚模复形，
由《叠的上同调》命题
\ref{stacks-cohomology-proposition-loc-qcoh-flat-base-change}
可知 $H^p(Lg_!\mathcal{H})$ 是拟凝聚的。
\end{proof}




\section{上同调与光滑-\'etale、平坦-fppf site}
\label{stacks-perfect-section-compare-unbounded-lisse-etale}

\noindent
我们已经看到，代数叠 $\mathcal{X}$ 上的层的上同调可在平坦-fppf site
% P12-ERR-0041: see ../control/ERRATA_CANDIDATES.jsonl
上计算。本节证明，对 $\mathcal{X}$ 的直接范畴中（可能）无界的对象，
同样的结论成立。

\begin{lemma}
\label{stacks-perfect-lemma-higher-shriek-Z}
设 $\mathcal{X}$ 为代数叠。无论 $Lg_!$ 取引理
\ref{stacks-perfect-lemma-shriek-derived} 第 (1) 部分中的函子，还是 $Lg_!$ 取引理
\ref{stacks-perfect-lemma-shriek-derived} 第 (3) 部分中的函子，
都有 $Lg_!\mathbf{Z} = \mathbf{Z}$。
\end{lemma}

\begin{proof}
我们先对平坦-fppf site 与 fppf site 的比较作证明；
光滑-\'etale site 的情形完全相同。
我们需要证明：当 $i \not = 0$ 时 $H^i(Lg_!\mathbf{Z})$ 为 $0$，
并且典范映射 $H^0(Lg_!\mathbf{Z}) \to \mathbf{Z}$ 是同构。
设 $f : \mathcal{U} \to \mathcal{X}$ 为满的平坦态射，
其中 $\mathcal{U}$ 是概形，并且 $f$ 还是局部有限表现的。
（例如，选择一个表示 $U \to \mathcal{X}$，并令
$\mathcal{U}$ 为与 $U$ 对应的代数叠。）
由《叠上的层》引理
\ref{stacks-sheaves-lemma-check-exactness-covering} 和
\ref{stacks-sheaves-lemma-surjective-flat-locally-finite-presentation}，
只需证明：当 $i \not = 0$ 时，拉回
$f^{-1}H^i(Lg_!\mathbf{Z})$ 为 $0$，并且
$H^0(Lg_!\mathbf{Z}) \to f^{-1}\mathbf{Z}$ 的拉回是同构。
由引理 \ref{stacks-perfect-lemma-lisse-etale-functorial-derived} 可得
$f^{-1}Lg_!\mathbf{Z} = L(g')_!\mathbf{Z}$，其中
$g' : \Sh(\mathcal{U}_{flat, fppf}) \to \Sh(\mathcal{U}_{fppf})$
是 $\mathcal{U}$ 对应的比较态射。
这将我们约化到下一段所研究的情形。

\medskip\noindent
设对某个概形 $X$ 有 $\mathcal{X} = (\Sch/X)_{fppf}$。
此时范畴 $\mathcal{X}_{flat, fppf}$ 有终对象 $e$，即 $X/X$；
而且函子
$u : \mathcal{X}_{flat, fppf} \to \mathcal{X}_{fppf}$
将 $e$ 映到终对象。由于 $\mathbf{Z}$ 是终对象上的自由阿贝尔层
（只要终对象存在），由《site 上的上同调》引理
\ref{sites-cohomology-lemma-existence-derived-lower-shriek}
中 $Lg_!$ 的构造本身可得 $Lg_!\mathbf{Z} = \mathbf{Z}$。
\end{proof}

\begin{lemma}
\label{stacks-perfect-lemma-lisse-etale-cohomology}
设 $\mathcal{X}$ 为代数叠。采用引理
\ref{stacks-perfect-lemma-shriek-derived} 中的记号。
\begin{enumerate}
\item 对 $D(\mathcal{X}_\etale)$ 中的 $K$，有
\begin{enumerate}
\item $R\Gamma(\mathcal{X}_\etale, K) =
R\Gamma(\mathcal{X}_{lisse,\etale}, g^{-1}K)$，以及
\item 对 $\mathcal{X}_{lisse,\etale}$ 的任意对象 $x$，有
$R\Gamma(x, K) =
R\Gamma(\mathcal{X}_{lisse,\etale}/x, g^{-1}K)$。
\end{enumerate}
\item 对 $D(\mathcal{X}_{fppf})$ 中的 $K$，有
\begin{enumerate}
\item $R\Gamma(\mathcal{X}_{fppf}, K) =
R\Gamma(\mathcal{X}_{flat,fppf}, g^{-1}K)$，以及
\item
% P12-ERR-0043: see ../control/ERRATA_CANDIDATES.jsonl
对 $\mathcal{X}_{flat,fppf}$ 的任意对象 $x$，有
$H^p(x, K) =
R\Gamma(\mathcal{X}_{flat,fppf}/x, g^{-1}K)$。
\end{enumerate}
\end{enumerate}
两种情形下，对模也有同样结论，因为 $g^{-1} = g^*$，
并且由《site 上的上同调》引理
\ref{sites-cohomology-lemma-modules-abelian-unbounded}，
上同调的计算没有区别。
\end{lemma}

\begin{proof}
我们先对平坦-fppf site 与 fppf site 的比较作证明；
光滑-\'etale site 的情形完全相同。
由引理 \ref{stacks-perfect-lemma-higher-shriek-Z}，有
$Lg_!\mathbf{Z} = \mathbf{Z}$。于是得到
\begin{align*}
R\Gamma(\mathcal{X}_{fppf}, K)
& =
R\Hom(\mathbf{Z}, K) \\
& =
R\Hom(Lg_!\mathbf{Z}, K) \\
& =
R\Hom(\mathbf{Z}, g^{-1}K) \\
& =
% P12-ERR-0044: see ../control/ERRATA_CANDIDATES.jsonl
R\Gamma(\mathcal{X}_{lisse,\etale}, g^{-1}K)
\end{align*}
% P12-ERR-0045: see ../control/ERRATA_CANDIDATES.jsonl
这证明了 (1)(a)。
(1)(b) 由 (1)(a) 推出。事实上，若 $x$ 位于概形 $U$ 之上，
则 site $\mathcal{X}_\etale/x$ 等价于
$(\Sch/U)_\etale$，而
% P12-ERR-0046: see ../control/ERRATA_CANDIDATES.jsonl
$\mathcal{X}_{lisse,\etale}$ 等价于 $U_{lisse, \etale}$。
\end{proof}
\section{拟凝聚模的导出范畴}
\label{stacks-perfect-section-derived}

\noindent
设 $\mathcal{X}$ 为代数叠。由于包含函子
$\QCoh(\mathcal{O}_\mathcal{X}) \to
\textit{Mod}(\mathcal{O}_\mathcal{X})$ 不是正合的，
我们不能把 $D_\QCoh(\mathcal{O}_\mathcal{X})$ 定义为
$D(\mathcal{O}_\mathcal{X})$ 中由上同调层拟凝聚的复形组成的充满子范畴。
作为替代，我们仿照《叠的上同调》评注
\ref{stacks-cohomology-remark-bousfield-colocalization}，
将拟凝聚模的导出范畴定义为一个商。

\medskip\noindent
回忆 $\textit{LQCoh}^{fbc}(\mathcal{O}_\mathcal{X}) \subset
\textit{Mod}(\mathcal{O}_\mathcal{X})$ 表示由具有平坦基变换性质的
局部拟凝聚 $\mathcal{O}_\mathcal{X}$-模组成的充满子范畴；
参见《叠的上同调》第
\ref{stacks-cohomology-section-loc-qcoh-flat-base-change} 节。
我们将缩写
$$
D_{\textit{LQCoh}^{fbc}}(\mathcal{O}_\mathcal{X}) =
D_{\textit{LQCoh}^{fbc}(\mathcal{O}_\mathcal{X})}(\mathcal{O}_\mathcal{X})
$$
由《导出范畴》引理
\ref{derived-lemma-cohomology-in-serre-subcategory} 和
《叠的上同调》命题
\ref{stacks-cohomology-proposition-loc-qcoh-flat-base-change} 第 (2) 部分，
可知 $D_{\textit{LQCoh}^{fbc}}(\mathcal{O}_\mathcal{X})$
是 $D(\mathcal{O}_\mathcal{X})$ 的严格充满饱和三角子范畴。

\medskip\noindent
令 $\textit{Parasitic}(\mathcal{O}_\mathcal{X}) \subset
\textit{Mod}(\mathcal{O}_\mathcal{X})$ 表示由寄生
$\mathcal{O}_\mathcal{X}$-模组成的充满子范畴；
参见《叠的上同调》第
\ref{stacks-cohomology-section-parasitic} 节。
我们将缩写
$$
D_{\textit{Parasitic}}(\mathcal{O}_\mathcal{X}) =
D_{\textit{Parasitic}(\mathcal{O}_\mathcal{X})}(\mathcal{O}_\mathcal{X})
$$
与前面一样，这是 $D(\mathcal{O}_\mathcal{X})$
的严格充满饱和三角子范畴，因为
$\textit{Parasitic}(\mathcal{O}_\mathcal{X})$
是 $\textit{Mod}(\mathcal{O}_\mathcal{X})$ 的 Serre 子范畴；
参见《叠的上同调》引理
\ref{stacks-cohomology-lemma-parasitic}。

\medskip\noindent
$\textit{Mod}(\mathcal{O}_\mathcal{X})$ 的弱 Serre 子范畴
$\textit{Parasitic}(\mathcal{O}_\mathcal{X}) \cap
\textit{LQCoh}^{fbc}(\mathcal{O}_\mathcal{X})$
的交仍是弱 Serre 子范畴。
类似地，我们缩写
\begin{align*}
D_{\textit{Parasitic} \cap \textit{LQCoh}^{fbc}}(\mathcal{O}_\mathcal{X})
& =
D_{\textit{Parasitic}(\mathcal{O}_\mathcal{X}) \cap
\textit{LQCoh}^{fbc}(\mathcal{O}_\mathcal{X})}(\mathcal{O}_\mathcal{X}) \\
& =
D_{\textit{Parasitic}}(\mathcal{O}_\mathcal{X})
\cap
D_{\textit{LQCoh}^{fbc}}(\mathcal{O}_\mathcal{X})
\end{align*}
与前面一样，这是 $D(\mathcal{O}_\mathcal{X})$
的严格充满饱和三角子范畴。因而更是
$D_{\textit{LQCoh}^{fbc}}(\mathcal{O}_\mathcal{X})$
的严格充满饱和三角子范畴。

\begin{definition}
\label{stacks-perfect-definition-derived}
设 $\mathcal{X}$ 为代数叠。采用上述记号，我们把
{\it 上同调层拟凝聚的 $\mathcal{O}_\mathcal{X}$-模导出范畴}
定义为 Verdier 商\footnote{此定义不同于文献中的定义；
参见 \cite[6.3]{olsson_sheaves}，但由引理
\ref{stacks-perfect-lemma-derived-quasi-coherent}，它与该定义一致。}
$$
D_\QCoh(\mathcal{O}_\mathcal{X}) =
D_{\textit{LQCoh}^{fbc}}(\mathcal{O}_\mathcal{X})/
D_{\textit{Parasitic} \cap \textit{LQCoh}^{fbc}}(\mathcal{O}_\mathcal{X})
$$
\end{definition}

\noindent
Verdier 商定义见《导出范畴》第
\ref{derived-section-quotients} 节。
$D_{\textit{LQCoh}^{fbc}}(\mathcal{O}_\mathcal{X})$
中的态射 $a : E \to E'$ 在
$D_\QCoh(\mathcal{O}_\mathcal{X})$ 中成为同构，当且仅当锥
$C(a)$ 的上同调层寄生；参见《导出范畴》引理
\ref{derived-lemma-operations}。

\medskip\noindent
考虑函子
$$
D_{\textit{LQCoh}^{fbc}}(\mathcal{O}_\mathcal{X})
\xrightarrow{H^i}
\textit{LQCoh}^{fbc}(\mathcal{O}_\mathcal{X})
\xrightarrow{Q}
\QCoh(\mathcal{O}_\mathcal{X})
$$
注意 $Q$ 零化子范畴
$\textit{Parasitic}(\mathcal{O}_\mathcal{X}) \cap
\textit{LQCoh}^{fbc}(\mathcal{O}_\mathcal{X})$；
参见《叠的上同调》引理
\ref{stacks-cohomology-lemma-adjoint-kernel-parasitic}。
由《导出范畴》引理
\ref{derived-lemma-universal-property-quotient}，
我们得到上同调函子
\begin{equation}
\label{stacks-perfect-equation-Hi-quasi-coherent}
H^i :
D_\QCoh(\mathcal{O}_\mathcal{X})
\longrightarrow
\QCoh(\mathcal{O}_\mathcal{X})
\end{equation}
此外，$E \in D_\QCoh(\mathcal{O}_\mathcal{X})$ 为零，
当且仅当对所有 $i \in \mathbf{Z}$ 都有 $H^i(E) = 0$，
因为由《叠的上同调》引理
\ref{stacks-cohomology-lemma-adjoint-kernel-parasitic}，
$Q$ 的核恰好等于
$\textit{Parasitic}(\mathcal{O}_\mathcal{X}) \cap
\textit{LQCoh}^{fbc}(\mathcal{O}_\mathcal{X})$。

\medskip\noindent
注意范畴
$\textit{Parasitic}(\mathcal{O}_\mathcal{X}) \cap
\textit{LQCoh}^{fbc}(\mathcal{O}_\mathcal{X})$ 和
$\textit{LQCoh}^{fbc}(\mathcal{O}_\mathcal{X})$
也是 \'etale 拓扑中模的阿贝尔范畴
$\textit{Mod}(\mathcal{X}_\etale, \mathcal{O}_\mathcal{X})$
的弱 Serre 子范畴；参见《叠的上同调》命题
\ref{stacks-cohomology-proposition-loc-qcoh-flat-base-change} 和引理
\ref{stacks-cohomology-lemma-parasitic}。
因此下述引理的陈述有意义。

\begin{lemma}
\label{stacks-perfect-lemma-compare-etale-fppf-QCoh}
设 $\mathcal{X}$ 为代数叠。缩写
$\mathcal{P}_\mathcal{X} = \textit{Parasitic}(\mathcal{O}_\mathcal{X}) \cap
\textit{LQCoh}^{fbc}(\mathcal{O}_\mathcal{X})$。
比较态射
$\epsilon : \mathcal{X}_{fppf} \to \mathcal{X}_\etale$
诱导交换图
$$
\xymatrix{
D_{\textit{Parasitic} \cap \textit{LQCoh}^{fbc}}(\mathcal{O}_\mathcal{X})
\ar[r] &
D_{\textit{LQCoh}^{fbc}}(\mathcal{O}_\mathcal{X}) \ar[r] &
D(\mathcal{O}_\mathcal{X}) \\
D_{\mathcal{P}_\mathcal{X}}(\mathcal{X}_\etale, \mathcal{O}_\mathcal{X})
\ar[r] \ar[u]^{\epsilon^*} &
D_{\textit{LQCoh}^{fbc}(\mathcal{O}_\mathcal{X})}(
\mathcal{X}_\etale, \mathcal{O}_\mathcal{X})
\ar[r] \ar[u]^{\epsilon^*} &
D(\mathcal{X}_\etale, \mathcal{O}_\mathcal{X})
\ar[u]^{\epsilon^*}
}
$$
此外，左边两个竖直箭头都是三角范畴的等价，因而还得到等价
$$
D_{\textit{LQCoh}^{fbc}(\mathcal{O}_\mathcal{X})}
(\mathcal{X}_\etale, \mathcal{O}_\mathcal{X})
/
D_{\mathcal{P}_\mathcal{X}}(\mathcal{X}_\etale, \mathcal{O}_\mathcal{X})
\longrightarrow
D_\QCoh(\mathcal{O}_\mathcal{X})
$$
\end{lemma}

\begin{proof}
由于 $\epsilon^*$ 正合，显然可得到引理陈述中的图。
我们将把《site 上的上同调》引理
\ref{sites-cohomology-lemma-compare-topologies-derived-adequate-modules}
应用于下述情形，以证明中间的竖直箭头是等价：
$\mathcal{C} = \mathcal{X}$，
$\tau = fppf$，
$\tau' = \etale$，
$\mathcal{O} = \mathcal{O}_\mathcal{X}$，
$\mathcal{A} = \textit{LQCoh}^{fbc}(\mathcal{O}_\mathcal{X})$，
而 $\mathcal{B}$ 是 $\mathcal{X}$ 中位于仿射概形之上的对象的集合。
要说明该引理适用，需要检验条件 (1)、(2)、(3)、(4)。
条件 (1) 和 (2) 由上面的讨论显然成立
（明确地说，这由《叠的上同调》命题
\ref{stacks-cohomology-proposition-loc-qcoh-flat-base-change} 得出）。
条件 (3) 成立是因为每个概形都有由仿射开集组成的 Zariski 开覆盖。
条件 (4) 由《下降》引理
\ref{descent-lemma-quasi-coherent-and-flat-base-change} 得出。

\medskip\noindent
我们略去下述验证：范畴的等价
$\epsilon^* : 
D_{\textit{LQCoh}^{fbc}(\mathcal{O}_\mathcal{X})}(
\mathcal{X}_\etale, \mathcal{O}_\mathcal{X})
\to
D_{\textit{LQCoh}^{fbc}}(\mathcal{O}_\mathcal{X})$
在上同调层寄生的复形子范畴之间诱导等价。
\end{proof}

\noindent
设 $\mathcal{X}$ 为代数叠。
由《叠的上同调》引理
\ref{stacks-cohomology-lemma-quasi-coherent-weak-serre}，
拟凝聚模范畴
$\QCoh(\mathcal{O}_{\mathcal{X}_{lisse,\etale}})$
构成
$\textit{Mod}(\mathcal{O}_{\mathcal{X}_{lisse,\etale}})$
的弱 Serre 子范畴，而
$\QCoh(\mathcal{O}_{\mathcal{X}_{flat,fppf}})$
构成
$\textit{Mod}(\mathcal{O}_{\mathcal{X}_{flat,fppf}})$
的弱 Serre 子范畴。因此可以考虑
$$
D_\QCoh(\mathcal{O}_{\mathcal{X}_{lisse,\etale}}) =
D_{\QCoh(\mathcal{O}_{\mathcal{X}_{lisse,\etale}})}(
\mathcal{O}_{\mathcal{X}_{lisse,\etale}})
\subset
D(\mathcal{O}_{\mathcal{X}_{lisse,\etale}})
$$
以及类似的
$$
D_\QCoh(\mathcal{O}_{\mathcal{X}_{flat,fppf}}) =
D_{\QCoh(\mathcal{O}_{\mathcal{X}_{flat,fppf}})}(
\mathcal{O}_{\mathcal{X}_{flat,fppf}})
\subset
D(\mathcal{O}_{\mathcal{X}_{flat,fppf}})
$$
与上面一样，它们都是严格充满饱和三角子范畴。
事实表明，$D_\QCoh(\mathcal{O}_\mathcal{X})$
与其中任意一个都等价。

\begin{lemma}
\label{stacks-perfect-lemma-derived-quasi-coherent}
设 $\mathcal{X}$ 为代数叠。令
$\mathcal{P}_\mathcal{X} = \textit{Parasitic}(\mathcal{O}_\mathcal{X}) \cap
\textit{LQCoh}^{fbc}(\mathcal{O}_\mathcal{X})$。
\begin{enumerate}
\item
设 $\mathcal{F}^\bullet$ 是
$D_{\textit{LQCoh}^{fbc}(\mathcal{O}_\mathcal{X})}
(\mathcal{X}_\etale, \mathcal{O}_\mathcal{X})$ 的对象。
对光滑-\'etale site，取《叠的上同调》引理
\ref{stacks-cohomology-lemma-lisse-etale} 中的 $g$，则
\begin{enumerate}
\item $g^*\mathcal{F}^\bullet$ 属于
$D_\QCoh(\mathcal{O}_{\mathcal{X}_{lisse,\etale}})$；
\item $g^*\mathcal{F}^\bullet = 0$，当且仅当
$\mathcal{F}^\bullet$ 属于
$D_{\mathcal{P}_\mathcal{X}}(\mathcal{X}_\etale, \mathcal{O}_\mathcal{X})$；
\item 对
$D_\QCoh(\mathcal{O}_{\mathcal{X}_{lisse,\etale}})$ 中的
$\mathcal{H}^\bullet$，
$Lg_!\mathcal{H}^\bullet$ 属于
$D_{\textit{LQCoh}^{fbc}(\mathcal{O}_\mathcal{X})}(
\mathcal{X}_\etale, \mathcal{O}_\mathcal{X})$；并且
\item 函子 $g^*$ 和 $Lg_!$ 给出互逆函子
$$
\xymatrix{
D_\QCoh(\mathcal{O}_\mathcal{X}) \ar@<1ex>[r]^-{g^*} &
D_\QCoh(\mathcal{O}_{\mathcal{X}_{lisse,\etale}})
\ar@<1ex>[l]^-{Lg_!}
}
$$
\end{enumerate}
\item
设 $\mathcal{F}^\bullet$ 是
$D_{\textit{LQCoh}^{fbc}}(\mathcal{O}_\mathcal{X})$ 的对象。
对平坦-fppf site，取《叠的上同调》引理
\ref{stacks-cohomology-lemma-lisse-etale} 中的 $g$，则
\begin{enumerate}
\item $g^*\mathcal{F}^\bullet$ 属于
$D_\QCoh(\mathcal{O}_{\mathcal{X}_{flat, fppf}})$；
\item $g^*\mathcal{F}^\bullet = 0$，当且仅当
$\mathcal{F}^\bullet$ 属于
$D_{\mathcal{P}_\mathcal{X}}(\mathcal{O}_\mathcal{X})$；
\item 对
$D_\QCoh(\mathcal{O}_{\mathcal{X}_{flat,fppf}})$ 中的
$\mathcal{H}^\bullet$，
$Lg_!\mathcal{H}^\bullet$ 属于
$D_{\textit{LQCoh}^{fbc}}(\mathcal{O}_\mathcal{X})$；并且
\item 函子 $g^*$ 和 $Lg_!$ 给出互逆函子
$$
\xymatrix{
D_\QCoh(\mathcal{O}_\mathcal{X}) \ar@<1ex>[r]^-{g^*} &
D_\QCoh(\mathcal{O}_{\mathcal{X}_{flat,fppf}}) \ar@<1ex>[l]^-{Lg_!}
}
$$
\end{enumerate}
\end{enumerate}
\end{lemma}

\begin{proof}
函子 $g^* = g^{-1}$ 正合，故 (1)(a)、(2)(a)、(1)(b) 和 (2)(b)
由《叠的上同调》引理
\ref{stacks-cohomology-lemma-quasi-coherent} 和
\ref{stacks-cohomology-lemma-parasitic-in-terms-flat-fppf} 得出。

\medskip\noindent
证明 (1)(c) 和 (2)(c)。
引理 \ref{stacks-perfect-lemma-shriek-derived} 中 $Lg_!$ 的构造
（通过《site 上的上同调》引理
\ref{sites-cohomology-lemma-existence-derived-lower-shriek}，
而该引理又使用《导出范畴》命题
\ref{derived-proposition-left-derived-exists}）
表明：对 $D(\mathcal{O}_{\mathcal{X}_{lisse,\etale}})$
中任意对象 $\mathcal{H}^\bullet$，$Lg_!$ 的计算式为
$$
Lg_!\mathcal{H}^\bullet = \colim g_!\mathcal{K}_n^\bullet =
g_! \colim \mathcal{K}_n^\bullet
$$
（逐项取余极限），其中拟同构
$\colim \mathcal{K}_n^\bullet \to \mathcal{H}^\bullet$
诱导拟同构
$\mathcal{K}_n^\bullet \to \tau_{\leq n} \mathcal{H}^\bullet$。
由于包含函子
$$
\textit{LQCoh}^{fbc}(\mathcal{O}_\mathcal{X}) \subset
\textit{Mod}(\mathcal{X}_\etale, \mathcal{O}_\mathcal{X})
\quad\text{且}\quad
\textit{LQCoh}^{fbc}(\mathcal{O}_\mathcal{X}) \subset
\textit{Mod}(\mathcal{O}_\mathcal{X})
$$
与滤过余极限相容，只需对
$D_\QCoh(\mathcal{O}_{\mathcal{X}_{lisse,\etale}})$ 和
$D_\QCoh(\mathcal{O}_{\mathcal{X}_{flat,fppf}})$ 中的上有界复形
$\mathcal{H}^\bullet$ 证明 (c)。
在此情形下，为证明 $H^n(Lg_!\mathcal{H}^\bullet)$ 属于
$\textit{LQCoh}^{fbc}(\mathcal{O}_\mathcal{X})$，
可对满足 $\mathcal{H}^i = 0$（当 $i > m$ 时）的整数 $m$ 作归纳。
若 $m < n$，则 $H^n(Lg_!\mathcal{H}^\bullet) = 0$，结论成立。
一般地，考虑优三角
$$
\tau_{\leq m - 1}\mathcal{H}^\bullet \to \mathcal{H}^\bullet \to
H^m(\mathcal{H}^\bullet)[-m] \to \ldots
$$
（《导出范畴》评注
\ref{derived-remark-truncation-distinguished-triangle}），
并作用函子 $Lg_!$。由于
$\textit{LQCoh}^{fbc}(\mathcal{O}_\mathcal{X})$
是模范畴的弱 Serre 子范畴，只需对三项中的两项证明 (c)。
对 $Lg_!\tau_{\leq m - 1}\mathcal{H}^\bullet$，
结论由归纳假设成立；对
$Lg_!H^m(\mathcal{H}^\bullet)[-m]$，
结论由引理 \ref{stacks-perfect-lemma-higher-shriek-quasi-coherent} 成立。
所以 (c) 成立。

\medskip\noindent
我们证明 (2)(d)。由 (2)(a) 和 (2)(b)，函子
$g^{-1} = g^*$ 诱导函子
$$
c :
D_\QCoh(\mathcal{O}_\mathcal{X})
\longrightarrow
D_\QCoh(\mathcal{O}_{\mathcal{X}_{flat, fppf}})
$$
参见《导出范畴》引理
\ref{derived-lemma-universal-property-quotient}。
因此有下述三角范畴的图
$$
\xymatrix{
D_{\textit{LQCoh}^{fbc}}(\mathcal{O}_\mathcal{X})
\ar[rd]^{g^{-1}} \ar[rr]_q & &
D_\QCoh(\mathcal{O}_\mathcal{X}) \ar[ld]^c \\
& D_\QCoh(\mathcal{O}_{\mathcal{X}_{flat, fppf}})
\ar@<1ex>[lu]^{Lg_!}
}
$$
其中 $q$ 为商函子，内三角交换，并且
$g^{-1}Lg_! = \text{id}$。
% P12-ERR-0047: see ../control/ERRATA_CANDIDATES.jsonl
对 $D_{\textit{LQCoh}^{fbc}}(\mathcal{O}_\mathcal{X})$
中 $E$ 的任意对象，映射 $a : Lg_!g^{-1}E \to E$
在 $D(\mathcal{O}_{\mathcal{X}_{flat, fppf}})$ 中映为拟同构。
因此 $a$ 的锥在 $g^{-1}$ 下映为零，由 (2)(b) 可知
$q(a)$ 是同构。所以 $q \circ Lg_!$ 是 $c$ 的拟逆。

\medskip\noindent
在光滑-\'etale site 的情形下，与上面完全相同的论证证明
$$
D_{\textit{LQCoh}^{fbc}(\mathcal{O}_\mathcal{X})}(
\mathcal{X}_\etale, \mathcal{O}_\mathcal{X})
/
D_{\mathcal{P}_\mathcal{X}}(
\mathcal{X}_\etale, \mathcal{O}_\mathcal{X})
$$
等价于
$D_\QCoh(\mathcal{O}_{\mathcal{X}_{lisse,\etale}})$。
应用引理 \ref{stacks-perfect-lemma-compare-etale-fppf-QCoh}
中的最后一个等价即可完成证明。
\end{proof}

\noindent
下述引理告诉我们，商函子
$D_{\textit{LQCoh}^{fbc}}(\mathcal{O}_\mathcal{X}) \to
D_\QCoh(\mathcal{O}_\mathcal{X})$
有左伴随。参见评注 \ref{stacks-perfect-remark-QCoh-admissible}。

\begin{lemma}
\label{stacks-perfect-lemma-bousfield-colocalization}
设 $\mathcal{X}$ 为代数叠。
设 $E$ 是 $D_{\textit{LQCoh}^{fbc}}(\mathcal{O}_\mathcal{X})$ 的对象。
在 $D_{\textit{LQCoh}^{fbc}}(\mathcal{O}_\mathcal{X})$ 中
存在典范优三角
$$
E' \to E \to P \to E'[1]
$$
使得 $P$ 属于
$D_{\textit{Parasitic} \cap \textit{LQCoh}^{fbc}}
(\mathcal{O}_\mathcal{X})$，并且对所有
$D_{\textit{Parasitic} \cap \textit{LQCoh}^{fbc}}(\mathcal{O}_\mathcal{X})$
中的 $P'$，都有
$$
\Hom_{D(\mathcal{O}_\mathcal{X})}(E', P') = 0
$$
\end{lemma}

\begin{proof}
考虑《叠的上同调》第
\ref{stacks-cohomology-section-lisse-etale} 节所研究的赋环拓扑斯态射
$g : \Sh(\mathcal{X}_{flat, fppf}) \longrightarrow \Sh(\mathcal{X}_{fppf})$。
令 $E' = Lg_!g^*E$，并令 $P$ 为伴随映射
$E' \to E$ 的锥；参见引理
\ref{stacks-perfect-lemma-shriek-derived} 第 (4) 部分。
由引理 \ref{stacks-perfect-lemma-derived-quasi-coherent} 的 (2)(a) 和 (2)(c)，
$E'$ 属于
$D_{\textit{LQCoh}^{fbc}}(\mathcal{O}_\mathcal{X})$。
所以 $P$ 也属于
$D_{\textit{LQCoh}^{fbc}}(\mathcal{O}_\mathcal{X})$。
由引理 \ref{stacks-perfect-lemma-shriek-derived} 第 (4) 部分，
$g^*Lg_! = \text{id}$，故映射 $g^*E' \to g^*E$ 是同构。
因此 $g^*P = 0$，进而由引理
\ref{stacks-perfect-lemma-derived-quasi-coherent} 的 (2)(b)，$P$ 是
$D_{\textit{Parasitic} \cap \textit{LQCoh}^{fbc}}(\mathcal{O}_\mathcal{X})$
的对象。
最后，对
$D_{\textit{Parasitic} \cap \textit{LQCoh}^{fbc}}(\mathcal{O}_\mathcal{X})$
中的 $P'$，有
$$
\Hom(E', P') = \Hom(Lg_!g^*E, P') = \Hom(g^*E, g^*P') = 0
$$
因为由引理 \ref{stacks-perfect-lemma-derived-quasi-coherent} 的 (2)(b)，
$g^*P' = 0$。
优三角 $E' \to E \to P \to E'[1]$ 是典范的
% P12-ERR-0048: see ../control/ERRATA_CANDIDATES.jsonl
（更精确地说，直到诱导 $E$ 上恒等映射的三角同构为止是唯一的）；
这由《导出范畴》第
\ref{derived-section-admissible} 节中的讨论得到。
\end{proof}

\begin{remark}
\label{stacks-perfect-remark-QCoh-admissible}
引理 \ref{stacks-perfect-lemma-bousfield-colocalization} 的结论告诉我们，
$$
D_{\textit{Parasitic} \cap \textit{LQCoh}^{fbc}}(\mathcal{O}_\mathcal{X})
\subset
D_{\textit{LQCoh}^{fbc}}(\mathcal{O}_\mathcal{X})
$$
是左可容许子范畴；参见《导出范畴》第
\ref{derived-section-admissible} 节。
特别地，若
$\mathcal{A} \subset D_{\textit{LQCoh}^{fbc}}(\mathcal{O}_\mathcal{X})$
表示它的左正交，则《导出范畴》命题
\ref{derived-proposition-summarize-admissible} 推出：
$\mathcal{A}$ 是
$D_{\textit{LQCoh}^{fbc}}(\mathcal{O}_\mathcal{X})$
中的右可容许子范畴，并且复合
$$
\mathcal{A} \longrightarrow
D_{\textit{LQCoh}^{fbc}}(\mathcal{O}_\mathcal{X}) \longrightarrow
D_\QCoh(\mathcal{O}_\mathcal{X})
$$
是等价。这意味着，可将
$D_\QCoh(\mathcal{O}_\mathcal{X})$
视为 $D_{\textit{LQCoh}^{fbc}}(\mathcal{O}_\mathcal{X})$
的严格充满饱和三角子范畴，也可视为
$D(\mathcal{X}_{fppf}, \mathcal{O}_\mathcal{X})$
的严格充满饱和三角子范畴。
\end{remark}





\section{拟凝聚模的导出直像}
\label{stacks-perfect-section-derived-pushforward}

\noindent
作为上面材料的第一个应用，我们构造导出直像。
读者可在《例》第
\ref{examples-section-derived-push-quasi-coherent} 节找到代数叠之间的
拟紧拟分离态射 $f : \mathcal{X} \to \mathcal{Y}$ 的一个例子，
其中直像函子 $Rf_*$ 并不诱导函子
$D_\QCoh(\mathcal{O}_\mathcal{X}) \to
D_\QCoh(\mathcal{O}_\mathcal{Y})$。
因此，有必要限制到下有界复形。

\begin{proposition}
\label{stacks-perfect-proposition-derived-direct-image-quasi-coherent}
设 $f : \mathcal{X} \to \mathcal{Y}$ 为代数叠之间的拟紧拟分离态射。
函子 $Rf_*$ 诱导交换图
$$
\xymatrix{
D^{+}_{\textit{Parasitic} \cap \textit{LQCoh}^{fbc}}(\mathcal{O}_\mathcal{X})
\ar[r] \ar[d]^{Rf_*} &
D^{+}_{\textit{LQCoh}^{fbc}}(\mathcal{O}_\mathcal{X})
\ar[r] \ar[d]^{Rf_*} &
D(\mathcal{O}_\mathcal{X})
\ar[d]^{Rf_*} \\
D^{+}_{\textit{Parasitic} \cap \textit{LQCoh}^{fbc}}(\mathcal{O}_\mathcal{Y})
\ar[r] &
D^{+}_{\textit{LQCoh}^{fbc}}(\mathcal{O}_\mathcal{Y}) \ar[r] &
D(\mathcal{O}_\mathcal{Y})
}
$$
从而在商范畴上诱导函子
$$
Rf_{\QCoh, *} :
D^{+}_\QCoh(\mathcal{O}_\mathcal{X})
\longrightarrow
D^{+}_\QCoh(\mathcal{O}_\mathcal{Y})
$$
此外，《叠的上同调》命题
\ref{stacks-cohomology-proposition-direct-image-quasi-coherent} 中的函子
% P12-ERR-0049: see ../control/ERRATA_CANDIDATES.jsonl
$R^if_\QCoh$ 等于 $H^i \circ Rf_{\QCoh, *}$，
其中 $H^i$ 如（\ref{stacks-perfect-equation-Hi-quasi-coherent}）所示。
\end{proposition}

\begin{proof}
需要证明：对
$D^{+}_{\textit{LQCoh}^{fbc}}(\mathcal{O}_\mathcal{X})$
中的 $E$，$Rf_*E$ 是
$D^{+}_{\textit{LQCoh}^{fbc}}(\mathcal{O}_\mathcal{Y})$ 的对象。
这由《叠的上同调》命题
\ref{stacks-cohomology-proposition-loc-qcoh-flat-base-change}
以及谱序列
$R^if_*H^j(E) \Rightarrow R^{i + j}f_*E$ 得出。
对寄生模的情形，用《叠的上同调》引理
\ref{stacks-cohomology-lemma-pushforward-parasitic}
作同样的论证。最后一个断言由
（\ref{stacks-perfect-equation-Hi-quasi-coherent}）中 $H^i$ 的定义显然成立。
\end{proof}




\section{拟凝聚模的导出拉回}
\label{stacks-perfect-section-derived-pullback}

\noindent
具有拟凝聚上同调层的复形的导出拉回总是存在。

\begin{proposition}
\label{stacks-perfect-proposition-derived-pullback-quasi-coherent}
设 $f : \mathcal{X} \to \mathcal{Y}$ 为代数叠之间的态射。
正合函子 $f^*$ 诱导交换图
$$
\xymatrix{
D_{\textit{LQCoh}^{fbc}}(\mathcal{O}_\mathcal{X}) \ar[r] &
D(\mathcal{O}_\mathcal{X}) \\
D_{\textit{LQCoh}^{fbc}}(\mathcal{O}_\mathcal{Y})
\ar[r] \ar[u]^{f^*} &
D(\mathcal{O}_\mathcal{Y}) \ar[u]^{f^*}
}
$$
复合
$$
D_{\textit{LQCoh}^{fbc}}(\mathcal{O}_\mathcal{Y})
\xrightarrow{f^*}
D_{\textit{LQCoh}^{fbc}}(\mathcal{O}_\mathcal{X})
\xrightarrow{q_\mathcal{X}}
D_\QCoh(\mathcal{O}_\mathcal{X})
$$
相对于局部化
$D_{\textit{LQCoh}^{fbc}}(\mathcal{O}_\mathcal{Y}) \to
D_\QCoh(\mathcal{O}_\mathcal{Y})$
是左可导的，并可把 $Lf^*_\QCoh$ 定义为它的左导出函子
$$
Lf_\QCoh^* :
D_\QCoh(\mathcal{O}_\mathcal{Y})
\longrightarrow
D_\QCoh(\mathcal{O}_\mathcal{X})
$$
（参见《导出范畴》定义
\ref{derived-definition-right-derived-functor-defined} 和
\ref{derived-definition-everywhere-defined}）。
若 $f$ 拟紧且拟分离，则 $Lf^*_\QCoh$ 和 $Rf_{\QCoh, *}$
满足下述伴随关系：
$$
\Hom_{D_\QCoh(\mathcal{O}_\mathcal{X})}(Lf^*_\QCoh A, B)
=
\Hom_{D_\QCoh(\mathcal{O}_\mathcal{Y})}(A, Rf_{\QCoh, *}B)
$$
其中 $A \in D_\QCoh(\mathcal{O}_\mathcal{Y})$ 且
$B \in D^{+}_\QCoh(\mathcal{O}_\mathcal{X})$。
\end{proposition}

\begin{proof}
为证明第一个断言，需要说明：对
$D_{\textit{LQCoh}^{fbc}}(\mathcal{O}_\mathcal{Y})$ 中的 $E$，
$f^*E$ 是
$D_{\textit{LQCoh}^{fbc}}(\mathcal{O}_\mathcal{X})$ 的对象。
由于 $f^* = f^{-1}$ 正合，而由《叠的上同调》命题
\ref{stacks-cohomology-proposition-loc-qcoh-flat-base-change}，
$f^*$ 将 $\textit{LQCoh}^{fbc}(\mathcal{O}_\mathcal{Y})$
映入 $\textit{LQCoh}^{fbc}(\mathcal{O}_\mathcal{X})$，
所以结论立即成立。

\medskip\noindent
令 $\mathcal{D} = D_{\textit{LQCoh}^{fbc}}(\mathcal{O}_\mathcal{Y})$。
令 $S$ 为 $\mathcal{D}$ 中由下述态射组成的类：
其锥是
$D_{\textit{Parasitic} \cap \textit{LQCoh}^{fbc}}(\mathcal{O}_\mathcal{Y})$
的对象。
令 $\mathcal{D}' = D_\QCoh(\mathcal{O}_\mathcal{X})$。
令 $F = q_\mathcal{X} \circ f^* : \mathcal{D} \to \mathcal{D}'$。
那么 $\mathcal{D}, S, \mathcal{D}', F$ 如《导出范畴》情形
\ref{derived-situation-derived-functor} 和定义
\ref{derived-definition-right-derived-functor-defined} 所述。
我们证明：对 $\mathcal{D}$ 的任意对象 $E$，$LF(E)$ 都有定义。
事实上，考虑引理 \ref{stacks-perfect-lemma-bousfield-colocalization} 中构造的三角
$$
E' \to E \to P \to E'[1]
$$
注意 $s : E' \to E$ 是 $S$ 的元素。我们断言 $E'$ 计算 $LF$。
事实上，设 $s' : E'' \to E$ 是 $S$ 的另一个元素，
也就是说，它嵌入三角
$E'' \to E \to P' \to E''[1]$，其中 $P'$ 属于
$D_{\textit{Parasitic} \cap \textit{LQCoh}^{fbc}}(\mathcal{O}_\mathcal{Y})$。
由引理 \ref{stacks-perfect-lemma-bousfield-colocalization}（及其证明），
$E' \to E$ 通过 $E'' \to E$ 分解。
因此 $E' \to E$ 在系统 $S/E$ 中是共终的。
于是显然 $E'$ 计算 $LF$。

\medskip\noindent
为证明最后一个断言，写
$B = q_\mathcal{X}(H)$ 和 $A = q_\mathcal{Y}(E)$。
如上选择 $E' \to E$。
一方面我们将使用
$Rf_{\QCoh, *}(B) = q_\mathcal{Y}(Rf_*H)$，
另一方面使用
$Lf^*_\QCoh(A) = q_\mathcal{X}(f^*E')$。
\begin{align*}
\Hom_{D_\QCoh(\mathcal{O}_\mathcal{X})}(Lf^*_\QCoh A, B)
& = 
\Hom_{D_\QCoh(\mathcal{O}_\mathcal{X})}(q_\mathcal{X}(f^*E'),
q_\mathcal{X}(H)) \\
& = 
\colim_{H \to H'} \Hom_{D(\mathcal{O}_\mathcal{X})}(f^*E', H') \\
& = \colim_{H \to H'} \Hom_{D(\mathcal{O}_\mathcal{Y})}(E', Rf_*H') \\
& = \Hom_{D(\mathcal{O}_\mathcal{Y})}(E', Rf_*H) \\
& =
\Hom_{D_\QCoh(\mathcal{O}_\mathcal{Y})}(A, Rf_{\QCoh, *}B)
\end{align*}
这里的余极限取遍
$D^+_{\textit{LQCoh}^{fbc}}(\mathcal{O}_\mathcal{X})$
中满足下述条件的态射 $s : H \to H'$：
其锥 $P(s)$ 是
$D^+_{\textit{Parasitic} \cap \textit{LQCoh}^{fbc}}(\mathcal{O}_\mathcal{X})$
的对象。
第一个等式已经在上面看到。
第二个等式由 Verdier 商的构造成立。
第三个等式由《site 上的上同调》引理
\ref{sites-cohomology-lemma-adjoint} 成立。
由命题 \ref{stacks-perfect-proposition-derived-direct-image-quasi-coherent}，
$Rf_*P(s)$ 是
$D^+_{\textit{Parasitic} \cap \textit{LQCoh}^{fbc}}(\mathcal{O}_\mathcal{Y})$
的对象，因而
$\Hom_{D(\mathcal{O}_\mathcal{Y})}(E', Rf_*P(s)) = 0$。
所以第四个等式成立。最后一个等式由 $E'$ 的构造成立。
\end{proof}








\section{导出范畴中的拟凝聚对象}
\label{stacks-perfect-section-QC}

\noindent
本节接续《叠上的层》第
\ref{stacks-sheaves-section-QC} 节。
设 $\mathcal{X}$ 为代数叠。在该节中，我们定义了三角范畴
$$
\mathit{QC}(\mathcal{X}) = \mathit{QC}(\mathcal{X}_{affine}, \mathcal{O})
$$
并证明：若 $\mathcal{X}$ 可由代数空间 $X$ 表示，则
$\mathit{QC}(\mathcal{X})$ 等价于
$D_\QCoh(\mathcal{O}_X)$。
事实表明，我们已经发展了恰好足够的理论来证明：
对任意代数叠，同一结论都成立。

\begin{lemma}
\label{stacks-perfect-lemma-cohomology-parasitic}
设 $\mathcal{X}$ 为代数叠。
设 $K$ 是 $D(\mathcal{X}_{fppf})$ 的对象，且其上同调层寄生。
那么，对 $\mathcal{X}$ 中位于概形 $U$ 之上的每个对象 $x$，
只要 $U \to \mathcal{X}$ 平坦，就有 $R\Gamma(x, K) = 0$。
\end{lemma}

\begin{proof}
以
$g : \Sh(\mathcal{X}_{flat, fppf}) \to \Sh(\mathcal{X}_{fppf})$
表示第 \ref{stacks-perfect-section-lisse-etale} 节讨论的拓扑斯态射。
设 $x$ 为 $\mathcal{X}$ 中位于概形 $U$ 之上的对象，
且 $U \to \mathcal{X}$ 平坦；也就是说，
$x$ 是 $\mathcal{X}_{flat, fppf}$ 的对象。
由引理 \ref{stacks-perfect-lemma-lisse-etale-cohomology} 的 (2)(b)，有
$R\Gamma(x, K) = R\Gamma(\mathcal{X}_{flat, fppf}/x, g^{-1}K)$。
然而，我们的假设意味着
$D(\mathcal{X}_{flat, fppf})$ 的对象 $g^{-1}K$ 的上同调层为零；
参见《叠的上同调》定义
\ref{stacks-cohomology-definition-parasitic}。
所以 $g^{-1}K = 0$，结论成立。
\end{proof}

\begin{lemma}
\label{stacks-perfect-lemma-cohomology-parasitic-complex}
设 $\mathcal{X}$ 为代数叠。
设 $K$ 是 $D(\mathcal{X}_{fppf})$ 的对象，并且对
$\mathcal{X}$ 中位于仿射概形 $U$ 之上的所有对象 $x$，
只要 $U \to \mathcal{X}$ 平坦，就有 $R\Gamma(x, K) = 0$。
那么对所有 $i$ 都有 $H^i(\mathcal{X}, K) = 0$。
\end{lemma}

\begin{proof}
以
$g : \Sh(\mathcal{X}_{flat, fppf}) \to \Sh(\mathcal{X}_{fppf})$
表示第 \ref{stacks-perfect-section-lisse-etale} 节讨论的拓扑斯态射。
由引理 \ref{stacks-perfect-lemma-lisse-etale-cohomology} 的 (2)(b)，
我们的假设意味着：在 $\mathcal{X}_{flat, fppf}$
中位于仿射概形之上的每个对象上，$g^{-1}K$ 的上同调都为零。
由于 $\mathcal{X}_{flat, fppf}$ 的每个对象 $x$ 都有由此类对象组成的覆盖，
可知 $g^{-1}K$ 的上同调层为零，即 $g^{-1}K = 0$。
于是当然有
$R\Gamma(\mathcal{X}_{flat, fppf}, g^{-1}K) = 0$，
再由引理 \ref{stacks-perfect-lemma-lisse-etale-cohomology} 的 (2)(a)，
这推出所需结论。
\end{proof}

\begin{lemma}
\label{stacks-perfect-lemma-higher-shriek-QC}
设 $\mathcal{X}$ 为代数叠。
设 $K$ 是
$D_\QCoh(\mathcal{O}_{\mathcal{X}_{flat, fppf}})$ 的对象。
那么 $Lg_!K$ 满足下述性质：对
$\mathcal{X}_{affine}$ 的任意态射 $x \to x'$，映射
$$
R\Gamma(x', Lg_!K) \otimes_{\mathcal{O}(x')}^\mathbf{L} \mathcal{O}(x)
\longrightarrow
R\Gamma(x, Lg_!K)
$$
是拟同构。
\end{lemma}

\begin{proof}
由引理 \ref{stacks-perfect-lemma-derived-quasi-coherent} 的 (2)(c)，对象
$Lg_!K$ 属于
$D_{\textit{LQCoh}^{fbc}}(\mathcal{O}_\mathcal{X})$。
由此很容易看出，若
$\mathcal{O}(x') \to \mathcal{O}(x)$ 是平坦环同态，
则引理中显示的映射是同构；细节从略。

\medskip\noindent
本段说明该问题对 \'etale 拓扑是局部的。
设 $x \to x'$ 是 $\mathcal{X}_{affine}$ 的一般态射。
设 $\{x'_i \to x'\}$ 是
$\mathcal{X}_{affine, \etale}$ 中的覆盖。
令 $x_i = x \times_{x'} x'_i$，则
$\{x_i \to x\}$ 也是
$\mathcal{X}_{affine, \etale}$ 的覆盖。
此时
$\mathcal{O}(x') \to \prod \mathcal{O}(x'_i)$
是忠实平坦 \'etale 环同态，并且
$$
\prod \mathcal{O}(x_i) =
\mathcal{O}(x) \otimes_{\mathcal{O}(x')} \left(\prod \mathcal{O}(x'_i)\right)
$$
因此，一个我们略去的简单代数论证表明：
只需对 $\mathcal{X}_{affine}$ 中的每个态射
$x_i \to x'_i$ 证明引理陈述中的结论。
换言之，该问题在 \'etale 拓扑中是局部的。

\medskip\noindent
选择概形 $X$ 以及满的光滑态射
$f : X \to \mathcal{X}$。
可以（按照我们的语言滥用）把 $f$ 视为 $\mathcal{X}$ 的对象；
于是 $(\Sch/X)_{fppf} = \mathcal{X}/f$，参见《叠上的层》第
\ref{stacks-sheaves-section-restriction} 节。
例如，由《叠上的层》引理
\ref{stacks-sheaves-lemma-surjective-flat-locally-finite-presentation}，
存在 \'etale 覆盖 $\{x'_i \to x'\}$，
使得 $x'_i : U'_i = p(x'_i) \to \mathcal{X}$ 通过 $f$ 分解。
由上一段的结果，可设 $x \to x'$ 是
$(\textit{Aff}/X)_{fppf}$ 中某个态射 $U \to U'$ 在函子
$(\Sch/X)_{fppf} \to \mathcal{X}$ 下的像。
% P12-ERR-0050: see ../control/ERRATA_CANDIDATES.jsonl
此时我们看出，$Lg_!K$ 在 $(\Sch/X)_{fppf}$ 上的限制等于
$f^*Lg_!K = L(g')_!(f')^*K$；这由引理
\ref{stacks-perfect-lemma-lisse-etale-functorial-derived} 得出。
这将我们约化到下一段讨论的情形。

\medskip\noindent
设 $\mathcal{X} = (\Sch/X)_{fppf}$，且 $x \to x'$
对应于仿射概形态射 $U \to U'$。
我们仍可在 $U'$ 上 \'etale（或 Zariski）局部地工作，
因而可设 $U' \to X$ 通过 $X$ 的某个仿射开集分解。
这将我们约化到下一段讨论的情形。

\medskip\noindent
设 $\mathcal{X} = (\Sch/X)_{fppf}$，其中
$X = \Spec(R)$ 是仿射概形，并且 $x \to x'$
对应于仿射概形态射 $U \to U'$。
令 $M^\bullet$ 为表示 $R\Gamma(X, K)$ 的 $R$-模复形。
由《代数进阶》引理
\ref{more-algebra-lemma-K-flat-resolution} 中的构造，
可设 $M^\bullet = \colim P_n^\bullet$，
其中每个 $P_n^\bullet$ 都是自由 $R$-模的上有界复形。
细节从略；另见《代数进阶》评注
\ref{more-algebra-remark-P-resolution}。
考虑 $X_{flat, fppf} = (\Sch/X)_{flat, fppf}$
上由规则
$$
U \longmapsto \Gamma(U, M^\bullet \otimes_R \mathcal{O}_U)
$$
给出的模复形 $M^\bullet_{flat, fppf}$。
由《下降》第
\ref{descent-section-quasi-coherent-sheaves} 节中的讨论，
这是层的复形。存在典范映射
$M^\bullet_{flat, fppf} \to K$；
由证明开头的评注，该映射在
$X_{flat, fppf}$ 的仿射对象上的截面上给出同构。
由于 $X_{flat, fppf}$ 的每个对象都有由仿射对象组成的覆盖，
可知 $M^\bullet_{flat, fppf}$ 与 $K$ 一致。

\medskip\noindent
令 $M^\bullet_{fppf}$ 为 $X_{fppf}$ 上由上述同一公式给出的模复形。
回忆 $Lg_!\mathcal{O} = g_!\mathcal{O} = \mathcal{O}$。
由于 $Lg_!$ 是 $g_!$ 的左导出函子，可得
$Lg_!P_{n, flat, fppf}^\bullet = P_{n, fppf}^\bullet$。
由于函子 $Lg_!$ 与同伦余极限交换
（也可使用它在《site 上的上同调》引理
\ref{sites-cohomology-lemma-existence-derived-lower-shriek} 中的构造），
并且 $M^\bullet = \colim P_n^\bullet$，可得
$Lg_!M^\bullet_{flat, fppf} = M^\bullet_{fppf}$。
设 $U = \Spec(A)$、$U' = \Spec(A')$，
并且 $U \to U'$ 对应于环同态 $A' \to A$。
由上可知
$$
R\Gamma(U, Lg_!K) = M^\bullet \otimes_R A
\quad\text{且}\quad
R\Gamma(U', Lg_!K) = M^\bullet \otimes_R A'
$$
由于 $M^\bullet$ 是 $R$-模的 K-平坦复形，
由张量积的传递性，
$$
R\Gamma(U', Lg_!K) \otimes_{A'}^\mathbf{L} A
\longrightarrow
R\Gamma(U, Lg_!K)
$$
是所需的拟同构。
\end{proof}

\begin{proposition}
\label{stacks-perfect-proposition-QC-compare}
设 $\mathcal{X}$ 为代数叠。那么
$\mathit{QC}(\mathcal{X})$ 典范等价于
$D_\QCoh(\mathcal{O}_\mathcal{X})$。
\end{proposition}

\begin{proof}
由《叠上的层》引理
\ref{stacks-sheaves-lemma-QC-compare-fppf}，
沿比较态射
$\epsilon : \mathcal{X}_{affine, fppf} \to \mathcal{X}_{affine}$
的拉回将 $\mathit{QC}(\mathcal{X})$ 与充满子范畴
$Q_\mathcal{X} \subset D(\mathcal{X}_{affine, fppf}, \mathcal{O})$
等同。利用《叠上的层》公式
（\ref{stacks-sheaves-equation-alternative-ringed}）中的赋环拓扑斯等价，
我们可以并将会把 $Q_\mathcal{X}$ 视为
$D(\mathcal{X}_{fppf}, \mathcal{O})$ 的充满子范畴。

\medskip\noindent
类似地，由引理 \ref{stacks-perfect-lemma-bousfield-colocalization} 和评注
\ref{stacks-perfect-remark-QCoh-admissible}，可将
$D_\QCoh(\mathcal{O}_\mathcal{X})$ 视为左可容许子范畴
$D_{\textit{Parasitic} \cap \textit{LQCoh}^{fbc}}(\mathcal{O}_\mathcal{X})$
在 $D_{\textit{LQCoh}^{fbc}}(\mathcal{O}_\mathcal{X})$
中的左正交 $\mathcal{A}$。

\medskip\noindent
为完成证明，我们将说明：作为
$D(\mathcal{X}_{fppf}, \mathcal{O})$ 的子范畴，
$Q_\mathcal{X}$ 等于 $\mathcal{A}$。

\medskip\noindent
步骤 1：$Q_\mathcal{X}$ 包含于
$D_{\textit{LQCoh}^{fbc}}(\mathcal{O}_\mathcal{X})$。
$Q_\mathcal{X}$ 的对象 $K$ 由下述性质刻画：
将 $K$ 视为 $D(\mathcal{X}_{affine, fppf}, \mathcal{O})$ 的对象时，
$R\epsilon_*K$ 是
$\mathit{QC}(\mathcal{X}_{affine}, \mathcal{O})$ 的对象。
这又恰好意味着：对 $\mathcal{X}_{affine}$ 的所有态射
$x \to x'$，映射
$$
R\Gamma(x', K) \otimes_{\mathcal{O}(x')}^\mathbf{L} \mathcal{O}(x)
\longrightarrow
R\Gamma(x, K)
$$
是同构；参见《site 上的上同调》引理
\ref{sites-cohomology-lemma-cartesion-plus-topology}
陈述中的脚注。
% P12-ERR-0051: see ../control/ERRATA_CANDIDATES.jsonl
现在，若 $x' \to x$ 位于仿射概形之间的某个平坦态射之上，
则这意味着
$$
H^i(x', K) \otimes_{\mathcal{O}(x')} \mathcal{O}(x)
\cong
H^i(x, K)
$$
这清楚地意味着 $H^i(K)$ 是 \'etale 拓扑的层
（《叠上的层》引理
\ref{stacks-sheaves-lemma-quasi-coherent-alternative}），
并且具有平坦基变换性质（略去一个小细节）。

\medskip\noindent
步骤 2：$Q_\mathcal{X}$ 包含于 $\mathcal{A}$。
为证明这一点，只需说明：对 $Q_\mathcal{X}$ 中的 $K$
以及
$D_{\textit{Parasitic} \cap \textit{LQCoh}^{fbc}}(\mathcal{O}_\mathcal{X})$
中的所有 $P$，都有 $\Hom(K, P) = 0$。
考虑对象
$$
H = R\SheafHom_{\mathcal{O}_\mathcal{X}}(K, P)
$$
设 $x$ 为 $\mathcal{X}$ 中位于仿射概形
$U = p(x)$ 之上的对象。
由《site 上的上同调》引理
\ref{sites-cohomology-lemma-section-RHom-over-U}，
下式中的第一个等式成立：
$$
R\Gamma(x, H) =
R\Hom_{\mathcal{O}_\mathcal{X}}(K|_{\mathcal{X}/x}, P|_{\mathcal{X}/x}) =
R\Hom_{\mathcal{O}}(K|_{\mathcal{X}_{affine}/x}, P|_{\mathcal{X}_{affine}/x})
$$
第二个等式来自下述事实：
site $\mathcal{X}/x$ 的拓扑斯等价于
site $\mathcal{X}_{affine}/x$ 的拓扑斯；
参见《叠上的层》公式
（\ref{stacks-sheaves-equation-alternative-ringed}）。
可对某个 $\mathit{QC}(\mathcal{O})$ 中的 $N$
写 $K = \epsilon^*N$。
于是由《site 上的上同调》引理
\ref{sites-cohomology-lemma-QC-hom-out-of-plus-topology}，有
$$
R\Gamma(x, H) =
R\Hom_{D(\mathcal{O}(x))}(R\Gamma(x, N), R\Gamma(x, P))
$$
由引理 \ref{stacks-perfect-lemma-cohomology-parasitic}，若
$U \to \mathcal{X}$ 平坦，则 $R\Gamma(x, P) = 0$，
因而在相同假设下 $R\Gamma(x, H) = 0$。
由引理 \ref{stacks-perfect-lemma-cohomology-parasitic-complex}，可得
$R\Gamma(\mathcal{X}, H) = 0$，从而 $\Hom(K, P) = 0$。

\medskip\noindent
步骤 3：$\mathcal{A}$ 包含于 $Q_\mathcal{X}$。
设 $K$ 是 $\mathcal{A}$ 的对象，并设
$x \to x'$ 为 $\mathcal{X}_{affine}$ 的态射。
需要证明
$$
R\Gamma(x', K) \otimes_{\mathcal{O}(x')}^\mathbf{L} \mathcal{O}(x)
\longrightarrow
R\Gamma(x, K)
$$
是拟同构；参见《site 上的上同调》引理
\ref{sites-cohomology-lemma-cartesion-plus-topology}
陈述中的脚注。
由引理 \ref{stacks-perfect-lemma-bousfield-colocalization} 的证明和评注
\ref{stacks-perfect-remark-QCoh-admissible} 中的讨论可知，
$\mathcal{A}$ 是 $Lg_!$ 在
$D_\QCoh(\mathcal{O}_{\mathcal{X}_{flat, fppf}})$
上的限制的像。因此可设对某个
$D_\QCoh(\mathcal{O}_{\mathcal{X}_{flat, fppf}})$ 中的 $M$，
有 $K = Lg_!M$。
% P12-ERR-0052: see ../control/ERRATA_CANDIDATES.jsonl
于是所需等式由引理 \ref{stacks-perfect-lemma-higher-shriek-QC} 得出。
\end{proof}
